(1) 轴对称Faraday定律 $\oint_{|\mathbf x|=r}\mathbf E\cdot d\mathbf l=-d\Phi_B/dt$。当 $r<a$，磁通为 $\pi r^2B(t)$，故 $\boxed{E_\phi(r,t)=-r\dot B(t)/2}$；当 $r>a$，通量为 $\pi a^2B(t)$，所以 $\boxed{E_\phi=-a^2\dot B(t)/(2r)}$，在 $r=a$ 连续。

(2) 取 $B=B_0\cos\omega t$，则 $E_\phi$ 与 $\sin\omega t$ 同相，内侧峰值 $r\omega B_0/2$、外侧峰值 $a^2\omega B_0/(2r)$。半径 $r$ 的回路电动势 $\mathcal E=-\pi\min(r^2,a^2)\dot B(t)$，峰值为 $\pi\min(r^2,a^2)\omega B_0$。

(3) 对 $r>a$ 的电阻环 $I(t)=\mathcal E(t)/R_c=\pi a^2\omega B_0\sin\omega t/R_c$。瞬时发热 $P=I^2R_c$，周期平均为 $\boxed{\bar P=\pi^2a^4\omega^2B_0^2/(2R_c)}$，全部来自维持变化磁通的外部能源。

(4) 因为 $\nabla\times\mathbf E=-\partial_t\mathbf B\ne0$，绕磁通的环流非零，不可能仅用全局单值标量势描述该电场。分开的开路端子测得电压还依赖测量线形成的回路及其链磁通；需要指定回路几何，而不是把感应电动势等同于静电势差。
